% Artículo VII. Recuperación por dependencia, 10-09-2026.
% RESULTADO_RECUPERADO: integral activo 2249, fuente/colaboracion/
% partes_iv_v/tex/residencias/86_cosmologia_cerrada.tex (S04),
% líneas 2--44 (cartas seleccionadas), 174--248 (tensor y aceleración).
% Fuente S04 leída íntegramente, SHA-256:
% feb2582c3316c8b30829cfb0476bfa2f4e95889dabfba16fe43f2d835a226673.
% Se recuperan definición, unicidad, conservación, descomposición,
% intercambio y aceleración, con sus condiciones explícitas.
% FORMALIZACION_REUNIDA: relación con la fuente efectiva de 31;
% se distingue conservación total de conservación material separada.
% CERTIFICADO_NUEVO: expansión de las pruebas de traza, balance y signos;
% cálculo local de Ricci para las tres cartas ya seleccionadas por S04.
% No se modifica el rebote de 50 ni se determina aquí una amplitud s_0.
%
% DEPENDENCIAS DE INTEGRACIÓN:
% - vii:sec:realizacion-cartan: la métrica procede de la cotetrada realizada.
% - vii:sec:acoplamientos: G_int y el reloj t_0 del mismo estado HMT;
%   c_int=1 es una elección de unidades posterior a su construcción.
% - vii:eq:ecuacion-metrica-efectiva y vii:eq:balance-metrico-total, de 31:
%   se usan únicamente para identificar el residual sobre una solución.
% - T_b es el tensor publicado por el transductor material declarado;
%   su conservación separada es una condición expresamente enunciada.
% - El selector TRIT de la carta y t_0(x)>0 son antecedentes de la
%   especialización de Sitter, no datos cosmológicos observacionales.
% El presente fragmento no certifica la clausura del conjunto de VII.

\section{Tenseur résiduel, décomposition de trace et bilan covariant}
\label{vii:sec:tensor-residual-balance}

La réalisation géométrique de l'état enrichi fournit une métrique, tandis que le transducteur matériel exprime la source à laquelle se rapporte la lecture. Le tenseur résiduel conserve la différence entre les deux représentations dans un même espace tensoriel. Sa décomposition permet de distinguer le secteur proportionnel à la métrique et le secteur de trace nulle, et détermine le courant d'échange correspondant à cette répartition.

\subsection{Données de la coordinatisation et construction du tenseur résiduel}
\label{vii:subsec:datos-tensor-residual}

Fixons un état HMT \(x\) et une coordinatisation connexe \(\mathcal U\) de sa réalisation géométrique. Sur \(\mathcal U\), la cotétrade non dégénérée de la section~\ref{vii:sec:realizacion-cartan} détermine une métrique lisse \(g=g_x\) de signature \((-+++)\). Nous notons \(\mathring\nabla\) sa connexion de Levi--Civita et écrivons
\[
 \mathsf G_{\mu\nu}[g]
 :=\operatorname{Ric}_{\mu\nu}[g]-\tfrac12R[g]g_{\mu\nu}.
\]
On distingue ainsi le tenseur d'Einstein \(\mathsf G[g]\) du couplage
gravitationnel \(G_{\rm int}\). La dérivée utilisée dans les bilans
de cette section est \(\mathring\nabla\), conformément à l'équation
métrique effective de la section~\ref{vii:subsec:fuente-metrica-efectiva}.

Nous travaillons dans la coordinatisation dimensionnelle \(c_{\rm int}=1\). Nous supposons \(G_{\rm int}(x)>0\) et \(\Lambda_{\rm ref}(x)\) constants sur \(\mathcal U\), et posons \(\kappa_x=8\pi G_{\rm int}(x)\). Le symbole \(x\) identifie l'état qui fournit ces lecteurs ; la différentiation covariante agit sur les points de \(\mathcal U\). La généalogie des couplages reste dans la section~\ref{v:ant:sec:gravity}.

Soit \(T_b\) le tenseur symétrique lisse exprimé par le transducteur matériel de la coordinatisation ; dans la lecture baryonique, il est sa source baryonique. Nous définissons
\begin{equation}
 \mathcal E^{\rm ref}_{\mu\nu}
 :=\kappa_x^{-1}\bigl(\mathsf G_{\mu\nu}[g]
                            +\Lambda_{\rm ref}g_{\mu\nu}\bigr),
 \qquad
 T^{\rm res}_{\mu\nu}:=\mathcal E^{\rm ref}_{\mu\nu}-T_{b,\mu\nu}.
 \label{vii:eq:residual-definicion}
\end{equation}
Le tenseur \(T^{\rm res}\) est le tenseur résiduel effectif, également noté \(T_{\rm osc}\) lorsqu'on interprète le secteur géométrique qui reste à exprimer par rapport à \(T_b\). La construction conserve comme arguments effectifs \(g_x,G_{\rm int},\Lambda_{\rm ref}\) et \(T_b\).

\begin{theorem}[Unicité du résiduel et bilan avec la source matérielle]
\label{vii:thm:residual-unicidad-balance}
Pour les données précédentes, il existe un unique tenseur symétrique \(T^{\rm res}\) qui satisfait \(\mathcal E^{\rm ref}=T_b+T^{\rm res}\). Sa divergence est
\begin{equation}
 \mathring\nabla^\mu T^{\rm res}_{\mu\nu}
 =-\mathring\nabla^\mu T_{b,\mu\nu}.
 \label{vii:eq:balance-residual-material}
\end{equation}
En particulier, si la source matérielle se conserve séparément,
\(\mathring\nabla^\mu T_{b,\mu\nu}=0\), \(T^{\rm res}\)
se conserve également.
\end{theorem}
\begin{proof}
La différence de deux tenseurs satisfaisant l'équation annoncée est nulle ; l'expression~\eqref{vii:eq:residual-definicion} fournit son existence. L'identité de Bianchi contractée et la compatibilité métrique de Levi--Civita donnent
\[
 \mathring\nabla^\mu\mathsf G_{\mu\nu}=0,
 \qquad \mathring\nabla g=0.
\]
Puisque \(\kappa_x\) et \(\Lambda_{\rm ref}\) sont constants sur la coordinatisation, \(\mathring\nabla^\mu\mathcal E^{\rm ref}_{\mu\nu}=0\). Dériver la différence qui définit \(T^{\rm res}\) démontre \eqref{vii:eq:balance-residual-material}, y compris l'affirmation sous conservation matérielle séparée.
\end{proof}

La constance spatiale et temporelle des couplages dans la coordinatisation fixe le domaine de ce bilan. Une collection d'états peut exprimer des couplages différents dans des coordinatisations différentes ; chaque identité précédente s'applique à la coordinatisation et à l'état qui la spécifient.

\subsection{Décomposition canonique et courant d'échange}
\label{vii:subsec:traza-corriente-residual}

Nous définissons la fonction de trace et les deux tenseurs
\begin{equation}
 \theta:=g^{\mu\nu}T^{\rm res}_{\mu\nu},\qquad
 T^{\Lambda}_{\mu\nu}:=\tfrac14\theta g_{\mu\nu},\qquad
 T^{0}_{\mu\nu}:=T^{\rm res}_{\mu\nu}-T^{\Lambda}_{\mu\nu}.
 \label{vii:eq:descomposicion-traza-residual}
\end{equation}
L'exposant \(\Lambda\) identifie la composante proportionnelle à la
métrique ; \(\Lambda_{\rm ref}\) est le paramètre de référence de
\eqref{vii:eq:residual-definicion}. Leurs fonctions sont distinctes.

\begin{proposition}[Projecteurs de trace et unicité de la décomposition]
\label{vii:prop:proyectores-traza-residual}
Sur les tenseurs symétriques covariants d'ordre deux, les applications
\[
 P_{\rm tr}(S)=\tfrac14\operatorname{tr}_g(S)g,
 \qquad P_0(S)=S-P_{\rm tr}(S)
\]
satisfont
\begin{equation}
 P_{\rm tr}^2=P_{\rm tr},\quad P_0^2=P_0,\quad
 P_{\rm tr}P_0=P_0P_{\rm tr}=0,\quad P_{\rm tr}+P_0=I.
 \label{vii:eq:proyectores-traza-residual}
\end{equation}
En particulier, \(T^{\rm res}=T^\Lambda+T^0\), \(\operatorname{tr}_gT^0=0\), et cette décomposition est unique.
\end{proposition}
\begin{proof}
En dimension quatre, \(g^{\mu\nu}g_{\mu\nu}=4\). Par conséquent,
\(\operatorname{tr}_g(P_{\rm tr}S)=\operatorname{tr}_gS\), ce qui
prouve l'idempotence de \(P_{\rm tr}\), l'annulation de la trace
de \(P_0S\) et, par développement, les autres identités.
Si \(S=fg+Q\) avec \(\operatorname{tr}_gQ=0\), prendre la trace
donne \(f=\operatorname{tr}_gS/4\). Cela détermine aussi
\(Q=S-fg\) et établit ainsi l'unicité.
\end{proof}

\begin{theorem}[Échange covariant entre les deux secteurs]
\label{vii:thm:intercambio-traza-residual}
Dans le domaine de \eqref{vii:eq:residual-definicion}, soit
\(b_\nu:=\mathring\nabla^\mu T_{b,\mu\nu}\). Le courant
d'échange de trace est la un-forme
\begin{equation}
 \mathcal J^{\rm tr}_\nu
 :=\tfrac14\partial_\nu\theta
 =\mathring\nabla^\mu T^\Lambda_{\mu\nu}.
 \label{vii:eq:corriente-intercambio-traza}
\end{equation}
Les bilans complets sont
\begin{equation}
 \mathring\nabla^\mu T^0_{\mu\nu}
       =-b_\nu-\mathcal J^{\rm tr}_\nu,
 \qquad
 \mathring\nabla^\mu(T^\Lambda+T^0)_{\mu\nu}=-b_\nu.
 \label{vii:eq:balances-sectores-residuales}
\end{equation}
Si \(b=0\), les deux secteurs échangent des courants exactement opposés. Dans ce cas, ils se conservent séparément si et seulement si \(\theta\) est constant dans la coordinatisation connexe.
\end{theorem}
\begin{proof}
La compatibilité métrique permet de calculer
\[
 \mathring\nabla^\mu\bigl(\tfrac14\theta g_{\mu\nu}\bigr)
 =\tfrac14g^{\mu\alpha}(\partial_\alpha\theta)g_{\mu\nu}
 =\tfrac14\partial_\nu\theta.
\]
Soustraire cette égalité à \eqref{vii:eq:balance-residual-material} produit
\eqref{vii:eq:balances-sectores-residuales}. Avec \(b=0\), la
conservation séparée équivaut à \(\mathrm d\theta=0\) ; une fonction
lisse de différentielle nulle est constante sur chaque composante connexe.
\end{proof}

La forme de degré un \(\mathcal J^{\rm tr}\) mesure l'échange de la répartition tensorielle. Elle conserve un type distinct de celui du covecteur de transport de la section~\ref{vii:sec:variacion-memoria} et du courant de spin \(\sigma_{IJ}\in\Omega^3(\mathcal U;\Lambda^2E^*)\) de \eqref{vii:eq:corriente-variacional}.

\subsection{Lien avec la source effective de Cartan--Holst}
\label{vii:subsec:residual-fuente-cartan}

Sur une solution de \eqref{vii:eq:ecuacion-metrica-efectiva}, dans les
mêmes unités \(c=1\), le résiduel a pour expression
\begin{equation}
 T^{\rm res}
 =T^{\rm phys,LC}+U^{\rm phys,spin}-T_b
   +\frac{\Lambda_{\rm ref}-\Lambda_{\rm int}}{\kappa_x}\,g.
 \label{vii:eq:residual-composicion-cartan}
\end{equation}
En effet, substituer \(\mathsf G+\Lambda_{\rm int}g
  =\kappa_x(T^{\rm phys,LC}+U^{\rm phys,spin})\) dans \eqref{vii:eq:residual-definicion} donne l'égalité terme à terme. Lorsque le transducteur matériel exprime précisément \(T_b=T^{\rm phys,LC}\) et \(\Lambda_{\rm ref}=\Lambda_{\rm int}\), on obtient \(T^{\rm res}=U^{\rm phys,spin}\). La condition de conservation séparée de \(T_b\) du théorème~\ref{vii:thm:residual-unicidad-balance} conserve sa fonction propre : \eqref{vii:eq:balance-metrico-total} établit le bilan de la source totale, et la répartition entre ses composantes dépend de l'action matérielle utilisée.

\subsection{Secteur de trace et accélération des coordinatisations sélectionnées}
\label{vii:subsec:residual-aceleracion-ds}

Le lecteur TRIT de l'état \(x\) exprime \(s=\tau(x)\in\{+1,0,-1\}\), et son horloge positive exprime
\begin{equation}
 H_*=\frac{2\pi}{108t_0(x)}>0.
 \label{vii:eq:reloj-carta-ds}
\end{equation}
Soit \(h_s\) une métrique tridimensionnelle de courbure sectionnelle constante \(s\). Les trois coordinatisations de la réalisation sélectionnée s'écrivent
\begin{equation}
 g_s=-\mathrm dt^2+a_s(t)^2h_s,\qquad
 a_+(t)=H_*^{-1}\cosh(H_*t),\quad
 a_0(t)=\exp(H_*t),\quad
 a_-(t)=H_*^{-1}\sinh(H_*t).
 \label{vii:eq:cartas-ds-residual}
\end{equation}
Les coordinatisations fermée et plate sont considérées pour \(t\in\mathbb R\) ; la coordinatisation ouverte, pour \(t>0\). Le signe \(s\) et l'horloge \(H_*\) maintiennent leurs antécédents dans l'état qui sélectionne la réalisation.

\begin{lemma}[Courbure des trois coordinatisations de la même horloge]
\label{vii:lem:curvatura-cartas-ds}
Pour les fonctions de \eqref{vii:eq:cartas-ds-residual},
\begin{equation}
 \frac{\ddot a_s}{a_s}=H_*^2,\qquad
 \left(\frac{\dot a_s}{a_s}\right)^2+\frac{s}{a_s^2}=H_*^2,
 \qquad
 \operatorname{Ric}[g_s]=3H_*^2g_s.
 \label{vii:eq:curvatura-cartas-ds}
\end{equation}
Par conséquent, \(R[g_s]=12H_*^2\) et
\(\mathsf G[g_s]=-3H_*^2g_s\).
\end{lemma}
\begin{proof}
Les deux premières égalités s'obtiennent en différentiant les trois fonctions : dans les coordinatisations fermée et ouverte, on utilise \(\cosh^2z-\sinh^2z=1\), et dans la coordinatisation plate \(\dot a_0=H_*a_0\). Pour vérifier l'identification tensorielle, les coefficients de connexion faisant intervenir la coordonnée temporelle sont
\[
 \Gamma^0{}_{ij}=a_s\dot a_s(h_s)_{ij},\qquad
 \Gamma^i{}_{0j}=(\dot a_s/a_s)\delta^i_j,
\]
et les composantes purement spatiales coïncident avec celles de \(h_s\). La contraction de la courbure donne
\[
 \operatorname{Ric}_{00}=-3\ddot a_s/a_s,\qquad
 \operatorname{Ric}_{0i}=0,\qquad
 \operatorname{Ric}_{ij}
  =(a_s\ddot a_s+2\dot a_s^2+2s)(h_s)_{ij}.
\]
Substituer les deux premières identités de
\eqref{vii:eq:curvatura-cartas-ds} produit
\(\operatorname{Ric}_{00}=-3H_*^2\) et
\(\operatorname{Ric}_{ij}=3H_*^2a_s^2(h_s)_{ij}\).
Prendre la trace et former le tenseur d'Einstein prouve les autres
affirmations, avec la convention de courbure de
\eqref{vii:eq:ecuacion-metrica-efectiva}.
\end{proof}

\begin{theorem}[Densité, pression et signe de l'accélération sélectionnée]
\label{vii:thm:residual-ds-aceleracion}
Dans les coordinatisations précédentes, pour \(\Lambda_{\rm ref}=0\) et \(T_b=0\), le tenseur résiduel est
\begin{equation}
 T^{\rm res}=-\rho_*g_s,\qquad
 \rho_*:=\frac{3H_*^2}{8\pi G_{\rm int}}>0,
 \qquad T^\Lambda=T^{\rm res},\quad T^0=0.
 \label{vii:eq:residual-ds-puro}
\end{equation}
Pour un observateur comobile unitaire \(u=\partial_t\), la densité
et la pression isotrope sont \(\rho_*\) et \(p_*=-\rho_*\).
L'équation d'accélération s'évalue sous la forme
\begin{equation}
 \frac{\ddot a_s}{a_s}
 =-\frac{4\pi G_{\rm int}}3(\rho_*+3p_*)=H_*^2>0.
 \label{vii:eq:aceleracion-residual-ds}
\end{equation}
\end{theorem}
\begin{proof}
Le lemme précédent et \eqref{vii:eq:residual-definicion} donnent \(T^{\rm res}=-(3H_*^2/\kappa_x)g_s\). Puisque \(g_s(u,u)=-1\), \(T^{\rm res}_{\mu\nu}u^\mu u^\nu=\rho_*\). Le projecteur spatial \(q^{\mu\nu}=g_s^{\mu\nu}+u^\mu u^\nu\) satisfait \(q^{\mu\nu}(g_s)_{\mu\nu}=3\), et donc
\[
 p_*:=\tfrac13q^{\mu\nu}T^{\rm res}_{\mu\nu}=-\rho_*.
\]
La trace vaut \(-4\rho_*\), de sorte que
\eqref{vii:eq:descomposicion-traza-residual} produit les deux secteurs
de \eqref{vii:eq:residual-ds-puro}. Enfin,
\(\rho_*+3p_*=-2\rho_*\), et la substitution dans le second membre
de \eqref{vii:eq:aceleracion-residual-ds} donne
\(\kappa_x\rho_*/3=H_*^2\), égal à la valeur géométrique déjà calculée.
La positivité utilise exactement \(G_{\rm int}>0\) et \(H_*>0\).
\end{proof}

Dans cette réalisation, le secteur de trace exprime une densité constante et une pression opposée, et l'horloge interne fixe l'accélération positive. Dans une coordinatisation générale, \(T^\Lambda\) décrit la composante isotrope ponctuelle, tandis que \(T^0\) conserve le résidu de trace nulle et que tous deux obéissent à l'échange de \eqref{vii:eq:balances-sectores-residuales}. L'identification de ces secteurs aux sources sombres observationnelles exige qu'une même réalisation reproduise conjointement l'expansion, les effets de lentille, la rotation, la croissance des structures et le fond cosmique. La présente construction fixe le tenseur, ses bilans et l'accélération des coordinatisations sélectionnées qui interviennent dans cette comparaison.
